\documentclass[10pt,a4paper]{amsart}
\usepackage{amsmath,amssymb,amsthm,graphicx}
% Retained minimal preamble for the Stacks Project chapter more-algebra.tex (More on Algebra).
% The chapter file pinned-original-source.tex carries only \input{preamble} on its first line, so no
% preamble text lives inside the chapter itself.  The Stacks Project uses one single shared preamble for
% every chapter; the snapshot stages it identically as sources/stacks-chow/preamble.tex and
% sources/stacks-etale-cohomology/preamble.tex.  The block below is copied mechanically, line for line,
% from sources/stacks-chow/preamble.tex lines 14-18 (the Xy-pic setup) and 154-217 (the native theorem
% declarations together with the mathematical macros).  Nothing was added, removed or redefined; the
% wrapper class and the amsmath/amssymb/amsthm/graphicx base packages are supplied by the assembler.
% Chapter: more-algebra.tex -- More on Algebra.
% For commutative diagrams we use Xy-pic
\usepackage[all]{xy}

% We use 2cell for 2-commutative diagrams.
\xyoption{2cell}

% Theorem environments.
%
\theoremstyle{plain}
\newtheorem{theorem}[subsection]{Theorem}
\newtheorem{proposition}[subsection]{Proposition}
\newtheorem{lemma}[subsection]{Lemma}

\theoremstyle{definition}
\newtheorem{definition}[subsection]{Definition}
\newtheorem{example}[subsection]{Example}
\newtheorem{exercise}[subsection]{Exercise}
\newtheorem{situation}[subsection]{Situation}

\theoremstyle{remark}
\newtheorem{remark}[subsection]{Remark}
\newtheorem{remarks}[subsection]{Remarks}

\numberwithin{equation}{subsection}

% Macros
%
\def\lim{\mathop{\mathrm{lim}}\nolimits}
\def\colim{\mathop{\mathrm{colim}}\nolimits}
\def\Spec{\mathop{\mathrm{Spec}}}
\def\Hom{\mathop{\mathrm{Hom}}\nolimits}
\def\Ext{\mathop{\mathrm{Ext}}\nolimits}
\def\SheafHom{\mathop{\mathcal{H}\!\mathit{om}}\nolimits}
\def\SheafExt{\mathop{\mathcal{E}\!\mathit{xt}}\nolimits}
\def\Sch{\mathit{Sch}}
\def\Mor{\mathop{\mathrm{Mor}}\nolimits}
\def\Ob{\mathop{\mathrm{Ob}}\nolimits}
\def\Sh{\mathop{\mathit{Sh}}\nolimits}
\def\NL{\mathop{N\!L}\nolimits}
\def\CH{\mathop{\mathrm{CH}}\nolimits}
\def\proetale{{pro\text{-}\acute{e}tale}}
\def\etale{{\acute{e}tale}}
\def\QCoh{\mathit{QCoh}}
\def\Ker{\mathop{\mathrm{Ker}}}
\def\Im{\mathop{\mathrm{Im}}}
\def\Coker{\mathop{\mathrm{Coker}}}
\def\Coim{\mathop{\mathrm{Coim}}}

% Boxtimes
%
\DeclareMathSymbol{\boxtimes}{\mathbin}{AMSa}{"02}

%
% Macros for moduli stacks/spaces
%
\def\QCohstack{\mathcal{QC}\!\mathit{oh}}
\def\Cohstack{\mathcal{C}\!\mathit{oh}}
\def\Spacesstack{\mathcal{S}\!\mathit{paces}}
\def\Quotfunctor{\mathrm{Quot}}
\def\Hilbfunctor{\mathrm{Hilb}}
\def\Curvesstack{\mathcal{C}\!\mathit{urves}}
\def\Polarizedstack{\mathcal{P}\!\mathit{olarized}}
\def\Complexesstack{\mathcal{C}\!\mathit{omplexes}}
% \Pic is the operator that assigns to X its picard group, usage \Pic(X)
% \Picardstack_{X/B} denotes the Picard stack of X over B
% \Picardfunctor_{X/B} denotes the Picard functor of X over B
\def\Pic{\mathop{\mathrm{Pic}}\nolimits}
\def\Picardstack{\mathcal{P}\!\mathit{ic}}
\def\Picardfunctor{\mathrm{Pic}}
\def\Deformationcategory{\mathcal{D}\!\mathit{ef}}

% --- appended 2026-10-12 (Chief): the project own definitions of the reference/slogan/history
% display environments, copied mechanically from the shared Stacks preamble lines 8-12
% (\usepackage{verbatim} plus the three \newenvironment lines, which the chapter preamble keeps
% verbatim from the project preamble). Nothing added, removed or redefined: these environments
% carry no mathematical content and the project itself typesets them as comments. Without these
% three lines a counted statement whose own text uses them fails with Environment slogan undefined.
\usepackage{verbatim}
\newenvironment{reference}{\comment}{\endcomment}
\newenvironment{slogan}{\comment}{\endcomment}
\newenvironment{history}{\comment}{\endcomment}
\usepackage{cleveref}
\crefname{originallemma}{original lemma}{original lemmas}
\Crefname{originallemma}{Original lemma}{Original lemmas}
\makeatletter
\newlabel{lemma-functorial}{{0624}{source}{}{sourcepointer-1}{}}
\newlabel{lemma-functorial@cref}{{[originallemma][0][]0624}{[1][0][]source}{}{}{}}
\makeatother
\crefname{originallemma}{original lemma}{original lemmas}
\Crefname{originallemma}{Original lemma}{Original lemmas}
\makeatletter
\newlabel{lemma-sequence-Koszul-complexes}{{0921}{source}{}{sourcepointer-2}{}}
\newlabel{lemma-sequence-Koszul-complexes@cref}{{[originallemma][0][]0921}{[1][0][]source}{}{}{}}
\makeatother
\crefname{originallemma}{original lemma}{original lemmas}
\Crefname{originallemma}{Original lemma}{Original lemmas}
\makeatletter
\newlabel{derived-lemma-pro-isomorphism}{{0G3C}{source}{}{sourcepointer-3}{}}
\newlabel{derived-lemma-pro-isomorphism@cref}{{[originallemma][0][]0G3C}{[1][0][]source}{}{}{}}
\makeatother
\crefname{originallemma}{original lemma}{original lemmas}
\Crefname{originallemma}{Original lemma}{Original lemmas}
\makeatletter
\newlabel{derived-lemma-essentially-constant-cohomology}{{0G3B}{source}{}{sourcepointer-4}{}}
\newlabel{derived-lemma-essentially-constant-cohomology@cref}{{[originallemma][0][]0G3B}{[1][0][]source}{}{}{}}
\makeatother
\title{The inverse system of powers of an ideal is pro-isomorphic to the system of Koszul complexes}
\author{The Stacks Project Authors}
\date{stacks-project:more-algebra.tex:local-snapshot-20261012; GNU FDL 1.2 or later}
\begin{document}
\maketitle
\setcounter{subsection}{1}
\noindent\textit{Editorial note.} All delivered mathematical text is the author's own, copied from the pinned Stacks Project chapter file more-algebra.tex; no statement or proof step was rewritten, reordered or removed. The counted statement is the labeled lemma environment at source lines 29647-29652 (canonical tag 0G3M) and the counted proof is the proof environment at source lines 29654-29674, which immediately follows it and is its own proof.\par\medskip
\noindent\textit{Editorial note.} The counted statement is formulated inside the situation of the section; that situation environment is delivered as the context span of this card, so the objects named in the statement can be read and the two labels occurring inside it resolve in the compiled document.\par\medskip
\noindent\textit{Editorial note.} Because the situation environment is delivered as the context span, no navigation stub is needed for the situation label; the stubs of this card exist for the three lemmas cited, and their kinds are the type words appearing in the labels themselves, not the names of the environments in which the labels live.\par\medskip
\noindent\textit{Editorial note.} The preamble span runs to line 90, the current end of the retained preamble, which includes the appended environment definitions; the counted text does not use them, and the span is taken whole so that the file version is unambiguous.\par\medskip
\noindent\textit{Editorial note.} The source URL is the source-level Stacks Project URL, as required when more than one card is produced from one source; the per-claim tag is carried by claim\_key and by this note. The counted statement and proof cite no bibliography entry, so this card needs no bibliography.\par\medskip
\noindent\textit{Editorial note.} The revision r002 corrects the boundaries of the context span: the situation environment opens at source line 26304 rather than 26305 and closes at line 26318, so the first assembly of revision r001 ended with a stray closing of that environment and was rejected; no mathematical text was changed.\par\medskip
\noindent\textit{Allowed standard prerequisites (editorial).}\begin{itemize}
\item The inverse system of powers of an ideal is pro-isomorphic to the system generated by the n-th powers of any finite generating set. used at proof lines 29654-29656. Replaces the system of powers by the system generated by the powers of the chosen generators, matching the situation.
\item If an inverse system of distinguished triangles has pro-zero cones then the two outer systems are pro-isomorphic. invoked at proof lines 29662-29664. Reduces the statement to showing that the system of cones is pro-zero.
\item For a bounded inverse system of complexes it suffices to check that the cohomology systems are pro-zero, and the vanishing of those systems was established in the earlier Koszul discussion using Noetherianity. used at proof lines 29664-29673. Supplies the pro-zero property of the cones and hence finishes the proof.
\end{itemize}\medskip
\ifdefined\hypertarget\hypertarget{sourcepointer-1}{}\fi
\noindent\textit{Source reference.} Original lemma 0624: sources/stacks-more-algebra/pinned-original-source.tex, source lines 6612--6612. The reference occurs inside the situation environment that is delivered as the context span, where it is used to produce the inverse system of Koszul complexes; it is a pointer needed only to define the objects of the statement.\par\medskip
\ifdefined\hypertarget\hypertarget{sourcepointer-2}{}\fi
\noindent\textit{Source reference.} Original lemma 0921: sources/stacks-more-algebra/pinned-original-source.tex, source lines 27172--27172. The reference is invoked at proof lines 29670-29672 for the pro-zero property of the Koszul cohomology systems, which is where the Noetherian hypothesis enters; it is the last input and everything leading to it is inside the counted span.\par\medskip
\ifdefined\hypertarget\hypertarget{sourcepointer-3}{}\fi
\noindent\textit{Source reference.} Original lemma 0G3C: sources/stacks-chow/tags.txt, source lines 20087--20087. The reference reduces the pro-isomorphism of the outer systems to pro-zero-ness of the cones, invoked at proof lines 29662-29664; the computations establishing that property are inside the counted span.\par\medskip
\ifdefined\hypertarget\hypertarget{sourcepointer-4}{}\fi
\noindent\textit{Source reference.} Original lemma 0G3B: sources/stacks-chow/tags.txt, source lines 20086--20086. The reference lets pro-zero-ness of a bounded system of complexes be checked on cohomology, invoked at proof lines 29664-29667; the degree computations that follow are inside the counted span.\par\medskip
\begin{situation}
\label{situation-koszul}
Let $A$ be a ring. Let $I = (f_1, \ldots, f_r) \subset A$. Let
$K_n^\bullet = K_\bullet(A, f_1^n, \ldots, f_r^n)$
be the Koszul complex on $f_1^n, \ldots, f_r^n$
viewed as a cochain complex in degrees $-r, -r + 1, \ldots, 0$.
Using the functoriality of Lemma \ref{lemma-functorial}
we obtain an inverse system
$$
\ldots \to K_3^\bullet \to K_2^\bullet \to K_1^\bullet
$$
compatible with the inverse system
$H^0(K_n^\bullet) = A/(f_1^n, \ldots, f_r^n)$
and compatible with the maps $A \to K_n^\bullet$.
\end{situation}

\begin{lemma}
\label{lemma-sequence-powers-pro-bounded}
In Situation \ref{situation-koszul} assume $A$ is Noetherian. With
notation as above, the inverse system $(I^n)$ is pro-isomorphic
in $D(A)$ to the inverse system $(I_n^\bullet)$.
\end{lemma}

\begin{proof}
It is elementary to show that the inverse system $I^n$ is pro-isomorphic
to the inverse system $(f_1^n, \ldots, f_r^n)$ in the category of $A$-modules.
Consider the inverse system of distinguished triangles
$$
I_n^\bullet \to (f_1^n, \ldots, f_r^n) \to C_n^\bullet \to I_n^\bullet[1]
$$
where $C_n^\bullet$ is the cone of the first arrow. By
Derived Categories, Lemma \ref{derived-lemma-pro-isomorphism}
it suffices to show that the inverse system $C_n^\bullet$ is pro-zero.
The complex $I_n^\bullet$ has nonzero terms only in degrees
$i$ with $-r + 1 \leq i \leq 0$ hence $C_n^\bullet$ is bounded
similarly. Thus by
Derived Categories, Lemma \ref{derived-lemma-essentially-constant-cohomology}
it suffices to show that $H^p(C_n^\bullet)$ is pro-zero. By the discussion
above we have $H^p(C_n^\bullet) = H^p(K_n^\bullet)$ for $p \leq -1$ and
$H^p(C_n^\bullet) = 0$ for $p \geq 0$.
The fact that the inverse systems $H^p(K_n^\bullet)$ are pro-zero
was shown in the proof of Lemma \ref{lemma-sequence-Koszul-complexes}
(and this is where the assumption that $A$ is Noetherian is used).
\end{proof}

\end{document}
